\RequirePackage{pdfmanagement}
\SetKeys[document/metadata]{pdfversion=1.7, pdfstandard=A-2u, lang=en-GB}
\input{preamble/pdfA-compliance/glyphtounicode}
\documentclass[11pt,a4paper]{report}

\newcommand*{\PaperTitle}{Sharp equilibrium bounds for \((n,n,n-1)\) mass-action networks: an answer to a question of Banaji and Feliu}
\newcommand*{\PaperTitlePlain}{Sharp equilibrium bounds for (n,n,n-1) mass-action networks: an answer to a question of Banaji and Feliu}
\newcommand*{\PaperShortTitle}{Sharp equilibrium bounds for \((n,n,n-1)\) networks}
\newcommand*{\PaperSubtitle}{chemistry-open-problems \textperiodcentered{} 18 October 2026 \textperiodcentered{} preliminary}
\newcommand*{\PaperCredits}{Solver: Claude Opus 5.5 (effort xhigh), orchestrating agent \textperiodcentered{} \(n=3\) case and computer searches: Claude Opus 5.5 exploration agents (effort high), 17 Oct 2026 \textperiodcentered{} Independent verification: pending}
\newcommand*{\PaperKeywords}{chemical reaction networks, multistationarity, mass-action kinetics, Descartes' rule of signs}
\input{chemop-preamble}
\addbibresource{references.bib}

\newcommand{\Net}{\mathcal N}
\newcommand{\one}{\mathbf 1}
\providecommand{\im}{\operatorname{im}}
\providecommand{\spn}{\operatorname{span}}
\providecommand{\rank}{\operatorname{rank}}

\begin{document}
\makepapertitle

\begin{paperabstract}
Banaji and Feliu proved that a mass-action network with \(n\) species, \(n\) reactions and rank \(n-1\) has at
most \(n\) positive nondegenerate equilibria on each stoichiometric class, and left open whether this bound is
attained for general \(n\) and, if so, what source molecularity is needed. We construct, for every \(n\ge3\), an
explicit \((n,n,n-1)\) network with sources \(X_1\) and \(2X_{i-1}+X_i\) that has exactly \(n\) positive
nondegenerate equilibria on one stoichiometric class for suitable rate constants. Together with their bound of two
equilibria for quadratic networks, this shows that the minimum source molecularity is exactly~3 for every
\(n\ge2\). The proof is short and computer-free: the reaction vectors form a path through the species, the
equilibria lie on a monomial curve, and the conservation law becomes a signomial with alternating coefficient signs.
\end{paperabstract}

\section*{The question}

Banaji and Feliu~\autocite{BanajiFeliu2026} study positive equilibria of mass-action networks with \(n\) species,
\(m\) reactions and stoichiometric rank \(r\), called \((n,m,r)\) networks. Their Corollary~4.21(i) states that
\((n,n,n-1)\) networks admit no more than \(n\) positive nondegenerate equilibria on any stoichiometric class. In
Remark~4.23 they write (arXiv:2409.06877v5, p.~30):

\begin{openquote}
``In the case \(n = 2\), both bounds in Corollary 4.21 are achieved by the \((2, 2, 1)\) network
\(X + 2Y \to 3Y\), \(Y \to X\) [\dots]. We remark, however, that quadratic \((2, 2, 1)\) networks do not admit
multiple positive nondegenerate equilibria [\dots]. We leave open whether the bound on positive nondegenerate
equilibria in Corollary 4.21(i) is achieved for general \(n\) and, if so, the minimum source molecularity of a
network needed to achieve the bound.''
\end{openquote}

We answer both parts.

\begin{theorem}\label{thm:A}
For every \(n\ge 2\) there is an \((n,n,n-1)\) mass-action network with source molecularity~3 which, for suitable
rate constants, has exactly \(n\) positive equilibria on some stoichiometric class, all of them nondegenerate.
Hence the bound of~\autocite[Corollary~4.21(i)]{BanajiFeliu2026} is attained for every \(n\).
\end{theorem}

\begin{corollary}\label{cor:B}
For every \(n\ge2\), the minimum source molecularity of an \((n,n,n-1)\) network attaining the bound \(n\) is
exactly~3.
\end{corollary}

\begin{proof}[of \zcref{cor:B} from \zcref{thm:A}]
For \(n=2\) this is the content of Remark~4.23 itself: the network \(X+2Y\to3Y,\ Y\to X\) has source molecularity~3,
and quadratic \((2,2,1)\) networks admit at most one positive nondegenerate equilibrium per class. For \(n\ge3\),
Corollary~5.16 of~\autocite{BanajiFeliu2026} states that quadratic \((n,n,n-1)\) networks (all sources of
molecularity at most~2) admit no more than two positive nondegenerate equilibria on each stoichiometric class, fewer
than \(n\). \zcref[S]{thm:A} provides networks of source molecularity~3.
\end{proof}

\section{Notation}

Species are \(X_1,\dots,X_n\) with concentrations \(x\in\mathbb R^n_{>0}\). Reaction \(R_j\) has source complex
\(a_j\in\mathbb Z^n_{\ge0}\), product complex \(b_j\in\mathbb Z^n_{\ge0}\) and reaction vector \(\gamma_j=b_j-a_j\). We
write \(\Gamma=[\gamma_1|\cdots|\gamma_m]\in\mathbb Z^{n\times m}\) for the stoichiometric matrix and
\(S\in\mathbb Z^{m\times n}\) for the matrix whose \(j\)-th row is \(a_j\). Mass-action kinetics with rate constants
\(\kappa\in\mathbb R^m_{>0}\) gives
\[
\dot x=\Gamma\, v(x),\qquad v_j(x)=\kappa_j\,x^{a_j}=\kappa_j\prod_i x_i^{a_{ji}},
\]
with Jacobian \(J(x)=\Gamma\,D_{v(x)}\,S\,D_{1/x}\), where \(D_u\) is the diagonal matrix with diagonal \(u\).
Following~\autocite[§2]{BanajiFeliu2026}, a positive equilibrium \(x\) is \emph{nondegenerate} if \(J(x)\) acts as a
nonsingular transformation on \(\im\Gamma\), i.e.\ \(\ker J(x)\cap\im\Gamma=\{0\}\). Stoichiometric classes are the
sets \((x_0+\im\Gamma)\cap\mathbb R^n_{\ge0}\). For \(\lambda>0\) and \(q\in\mathbb R^n\) we write
\(\lambda^q=(\lambda^{q_1},\dots,\lambda^{q_n})\); \(\circ\) is the componentwise product and \(e_i\) the \(i\)-th unit
vector.

\section{The family of networks}

Fix \(n\ge3\), define the signs
\[
\varepsilon_i=(-1)^{\lfloor (i-1)/2\rfloor}\quad(i=1,\dots,n),\qquad\text{i.e. } \varepsilon=(+,+,-,-,+,+,-,-,\dots),
\]
and \(s=-\varepsilon_n\in\{\pm1\}\). The network \(\Net_n\) consists of the \(n\) reactions
\[
\begin{aligned}
R_1&:\ X_1\ \longrightarrow\ X_1+\gamma_1,\\
R_{2k}&:\ 2X_{2k-1}+X_{2k}\ \longrightarrow\ 2X_{2k-1}+X_{2k}+\gamma_{2k}\qquad(2k\le n),\\
R_{2k+1}&:\ 2X_{2k}+X_{2k+1}\ \longrightarrow\ 2X_{2k}+X_{2k+1}+\gamma_{2k+1}\qquad(2k+1\le n),
\end{aligned}
\]
with reaction vectors
\[
\begin{aligned}
\gamma_1&=s\,e_1+e_n,\\
\gamma_{2k}&=c_k\,(e_{2k}-e_{2k-1}), & c_k&=(-1)^{k+1}s,\\
\gamma_{2k+1}&=d_k\,(e_{2k}+e_{2k+1}), & d_k&=(-1)^{k}s .
\end{aligned}
\]
Written out, the reactions are listed in \zcref{tab:reactions}.

\begin{table}[h]
\centering\small
\caption{The reactions of \(\Net_n\).}\label{tab:reactions}
\begin{tabular}{lll}
\toprule
reaction & coefficient \(+1\) & coefficient \(-1\)\\
\midrule
\(R_1\) (\(s\)) & \(X_1\to 2X_1+X_n\) & \(X_1\to X_n\)\\
\(R_{2k}\) (\(c_k\)) & \(2X_{2k-1}+X_{2k}\to X_{2k-1}+2X_{2k}\) & \(2X_{2k-1}+X_{2k}\to 3X_{2k-1}\)\\
\(R_{2k+1}\) (\(d_k\)) & \(2X_{2k}+X_{2k+1}\to 3X_{2k}+2X_{2k+1}\) & \(2X_{2k}+X_{2k+1}\to X_{2k}\)\\
\bottomrule
\end{tabular}
\end{table}

All products are nonnegative; the sources are \(a_1=e_1\) and \(a_j=2e_{j-1}+e_j\) (\(2\le j\le n\)), so the source
molecularity is~3 and the product molecularity is at most~5. For instance
\begin{align*}
\Net_3&:\ X_1\to 2X_1+X_3,\quad 2X_1+X_2\to X_1+2X_2,\quad 2X_2+X_3\to X_2;\\
\Net_4&:\ X_1\to X_4,\quad 2X_1+X_2\to 3X_1,\quad 2X_2+X_3\to 3X_2+2X_3,\quad 2X_3+X_4\to X_3+2X_4.
\end{align*}

\begin{lemma}[note=stoichiometry]\label{lem:stoich}
\(\Net_n\) is an \((n,n,n-1)\) network with \(\im\Gamma=\varepsilon^\perp\) (so \(\ker\Gamma^{\mathsf T}=\spn(\varepsilon)\))
and \(\ker\Gamma=\spn(\one)\), \(\one=(1,\dots,1)\). In particular it is dynamically nontrivial and its
stoichiometric classes are the sets \(\{x\ge0:\varepsilon\cdot x=T\}\), \(T\in\mathbb R\).
\end{lemma}

\begin{proof}
For \(t=1,\dots,n-1\) let \(u_t=e_{t+1}-e_t\) if \(t\) is odd and \(u_t=e_t+e_{t+1}\) if \(t\) is even. Then
\(\gamma_{t+1}\) is a nonzero multiple of \(u_t\) (\(\gamma_{2k}=c_ku_{2k-1}\), \(\gamma_{2k+1}=d_ku_{2k}\)). If \(t\) is
odd, \(\varepsilon_t=\varepsilon_{t+1}\); if \(t\) is even, \(\varepsilon_{t+1}=-\varepsilon_t\). Hence
\(\varepsilon\cdot u_t=0\) for all \(t\); also \(\varepsilon\cdot\gamma_1=s+\varepsilon_n=0\). The vectors
\(u_1,\dots,u_{n-1}\) are linearly independent: deleting the first coordinate leaves an upper bidiagonal
\((n-1)\times(n-1)\) matrix with unit diagonal (\(u_t\) is supported on \(\{t,t+1\}\) with coefficient \(1\) at
\(t+1\)). Therefore \(\im\Gamma=\spn(u_1,\dots,u_{n-1})=\varepsilon^\perp\) and \(\rank\Gamma=n-1\).

We check \(\sum_j\gamma_j=0\) coordinatewise. Coordinate~1: \(s-c_1=0\). Coordinate \(2k\): \(c_k+d_k=0\) if \(2k<n\);
if \(2k=n\) the coordinate is \(c_k+1\), and \(s=-\varepsilon_{2k}=-(-1)^{k-1}\) gives \(c_k=-1\). Coordinate
\(2k+1\ge3\): \(d_k-c_{k+1}=(-1)^ks-(-1)^{k+2}s=0\) if \(2k+1<n\); if \(2k+1=n\) the coordinate is \(d_k+1\), and
\(s=-\varepsilon_{2k+1}=-(-1)^k\) gives \(d_k=-1\). So \(\Gamma\one=0\), and since \(\rank\Gamma=n-1=m-1\),
\(\ker\Gamma=\spn(\one)\).
\end{proof}

\begin{lemma}[note=exponents]\label{lem:exponents}
The source matrix \(S\) is lower bidiagonal with unit diagonal, hence invertible, and \(q:=S^{-1}\one\) has entries
\[
q_i=\frac{1+2(-2)^{i-1}}{3}\qquad(i=1,\dots,n),\qquad q=(1,-1,3,-5,11,-21,43,\dots).
\]
These are pairwise distinct nonzero integers, ordered as
\(q_{2\lfloor n/2\rfloor}<\cdots<q_4<q_2<0<q_1<q_3<q_5<\cdots\).
\end{lemma}

\begin{proof}
\(Sq=\one\) reads \(q_1=1\) and \(2q_{i-1}+q_i=1\), i.e.\ \(q_i-\tfrac13=-2(q_{i-1}-\tfrac13)\). For odd \(i\),
\(q_i=(1+2^i)/3>0\) increases with \(i\); for even \(i\), \(q_i=(1-2^i)/3<0\) decreases with \(i\).
\end{proof}

\section{Equilibria lie on a monomial curve}

\begin{lemma}\label{lem:curve}
Fix \(\kappa\in\mathbb R^n_{>0}\) and put \(C=\exp(-S^{-1}\log\kappa)\in\mathbb R^n_{>0}\) (componentwise). The positive
equilibria of \(\Net_n\) are exactly the points \(x(\lambda)=C\circ\lambda^{q}\), \(\lambda>0\), and at \(x(\lambda)\)
every reaction has rate \(v_j=\lambda\). Conversely, every \(C\in\mathbb R^n_{>0}\) arises from the rate constants
\(\kappa_j=C^{-a_j}\).
\end{lemma}

\begin{proof}
By \zcref{lem:stoich}, \(x>0\) is an equilibrium iff \(v(x)\in\ker\Gamma=\spn(\one)\), i.e.\
\(\kappa_jx^{a_j}=\lambda\) for all \(j\) and some \(\lambda>0\). Taking logarithms,
\(S\log x=(\log\lambda)\one-\log\kappa\), so \(\log x=(\log\lambda)\,q+\log C\). The converse follows from
\(\log\kappa=-S\log C\).
\end{proof}

On the class \(\varepsilon\cdot x=T\), positive equilibria therefore correspond bijectively (\(\lambda\mapsto
x(\lambda)\) is injective because \(q\ne0\)) to the positive zeros of
\[
F(\lambda)=\sum_{i=1}^n\varepsilon_iC_i\,\lambda^{q_i}-T .
\]

\begin{lemma}[note=nondegeneracy]\label{lem:nondeg}
The equilibrium \(x(\lambda)\) is nondegenerate if and only if \(F'(\lambda)\neq0\).
\end{lemma}

\begin{proof}
At \(x=x(\lambda)\) we have \(v=\lambda\one\), so \(J=\lambda\,\Gamma SD_{1/x}\). Then \(Jy=0\) iff
\(SD_{1/x}y\in\ker\Gamma=\spn(\one)\) iff \(D_{1/x}y\in\spn(S^{-1}\one)=\spn(q)\) iff \(y\in\spn(x\circ q)\).
Since \(x'(\lambda)=\lambda^{-1}\,x(\lambda)\circ q\), we get \(\ker J=\spn(x'(\lambda))\). Because
\(J(\mathbb R^n)\subseteq\im\Gamma=\varepsilon^\perp\), \(\ker J\cap\im\Gamma=\{0\}\) iff
\(x'(\lambda)\notin\varepsilon^\perp\) iff \(\varepsilon\cdot x'(\lambda)=F'(\lambda)\neq0\).
\end{proof}

\section{Signomials with prescribed roots}

\begin{lemma}[note={Descartes' rule for real exponents}]\label{lem:descartes}
Let \(e_0<e_1<\dots<e_N\) be real and \(c\in\mathbb R^{N+1}\setminus\{0\}\). The number of zeros of
\(f(t)=\sum_kc_kt^{e_k}\) in \((0,\infty)\), counted with multiplicity, is at most the number of sign changes of
\((c_0,\dots,c_N)\) (zero entries ignored).
\end{lemma}

This is classical (Laguerre; see~\autocite[Part~V, Ch.~1]{PolyaSzego} or~\autocite{Jameson2006}).

\begin{proof}[note=sketch]
Induct on the number \(V\) of sign changes. If \(V=0\), \(f\) has no positive zero. Otherwise pick \(\theta\)
strictly between two consecutive exponents across which the sign changes, and consider
\(g(t)=(t^{-\theta}f(t))'=\sum_kc_k(e_k-\theta)t^{e_k-\theta-1}\). Multiplication by \(e_k-\theta\) flips exactly
the signs on one side of the chosen change, so \(g\) has \(V-1\) sign changes, while by Rolle's theorem \(g\) has at
least \(Z(f)-1\) positive zeros counted with multiplicity.
\end{proof}

\begin{lemma}\label{lem:prescribed}
Let \(e_0<\dots<e_N\) be real and \(0<\lambda_1<\dots<\lambda_N\). There is a nonzero \(c\in\mathbb R^{N+1}\) with
\(f(\lambda_r)=\sum_kc_k\lambda_r^{e_k}=0\) for all \(r\). Every such \(c\) has all entries nonzero and strictly
alternating in sign, and \(\lambda_1,\dots,\lambda_N\) are then all the positive zeros of \(f\), each of them
simple.
\end{lemma}

\begin{proof}
Existence: \(N\) homogeneous linear equations in \(N+1\) unknowns. By \zcref{lem:descartes}, \(f\) having at least
\(N\) positive zeros forces at least \(N\) sign changes in \((c_0,\dots,c_N)\); with \(N+1\) entries all of them are
nonzero and they alternate. \zcref[S]{lem:descartes} again shows that \(f\) has at most \(N\) positive zeros counted
with multiplicity, so the \(\lambda_r\) are all of them and each is simple.
\end{proof}

\section{Proof of the main theorem}

\begin{proof}[of \zcref{thm:A}]
The case \(n=2\) is in~\autocite[Remark~4.23]{BanajiFeliu2026}. Let \(n\ge3\), choose
\(0<\lambda_1<\dots<\lambda_n\), and list the \(n+1\) distinct exponents \(\{q_1,\dots,q_n,0\}\) in increasing order
(\zcref{lem:exponents}):
\(q_{2\lfloor n/2\rfloor}<\dots<q_4<q_2<0<q_1<q_3<\dots\).
\zcref[S]{lem:prescribed} gives a strictly alternating coefficient vector \(c\) whose signomial vanishes exactly at
the \(\lambda_r\), all simple; replacing \(c\) by \(-c\) if necessary, the coefficient of \(\lambda^{q_1}\) is
positive. Read the signs \(\varepsilon_i\) in the same order: \(\varepsilon_{2k}=(-1)^{k-1}\) alternates along the
negative exponents \(q_2,q_4,\dots\), \(\varepsilon_{2k+1}=(-1)^k\) alternates along the positive exponents
\(q_1,q_3,\dots\), and \(\varepsilon_1=\varepsilon_2=+1\). So the sequence of the \(\varepsilon\)'s, with an entry \(-1\)
inserted at exponent \(0\), alternates and starts at \(\lambda^{q_1}\) with \(+\), exactly like \(c\). Hence the
coefficient \(c_{(q_i)}\) of \(\lambda^{q_i}\) has sign \(\varepsilon_i\) and the constant coefficient \(c_{(0)}\) is
negative. Define
\[
C_i=\varepsilon_i\,c_{(q_i)}>0,\qquad T=-c_{(0)}>0,\qquad \kappa_j=C^{-a_j}.
\]
By \zcref{lem:curve}, the positive equilibria of \(\Net_n\) with these rate constants on the class
\(\varepsilon\cdot x=T\) correspond to the positive zeros of \(F=f\), namely \(\lambda_1,\dots,\lambda_n\). They are
simple, so the \(n\) equilibria \(x(\lambda_r)\) are nondegenerate by \zcref{lem:nondeg}, and there are no others on
that class.
\end{proof}

\begin{remark}
\begin{enumerate}
\item If the \(\lambda_r\) are rational, \(c\) can be chosen rational, so \(C\), \(T\), \(\kappa\) and the equilibria
\(x(\lambda_r)=C\circ\lambda_r^q\) are rational (the exponents are integers).
\item The stoichiometric classes of \(\Net_n\) are unbounded (\(\varepsilon\) has entries of both signs), as they must
be by~\autocite[Corollary~4.21(ii)]{BanajiFeliu2026} once \(n\ge3\).
\item \(\Net_n\) is an \((n,n,n-1)\) network with trivial partition, so it also shows that the bound
of~\autocite[Theorem~4.20(i)]{BanajiFeliu2026} is attained for every \(n\) when \(m=n\).
\item The exponents \(q_i\) grow like \(2^i\), which makes the rational witnesses large. Computer searches found other
networks with linearly growing exponents, but without a uniform proof.
\end{enumerate}
\end{remark}

\section{A small explicit instance}

For \(n=3\) take \(\lambda_r=r\). The exponents are \(q=(1,-1,3)\) and
\(\lambda F(\lambda)=-C_3\lambda^4+C_1\lambda^2-T\lambda+C_2\) must equal
\(-C_3(\lambda-1)(\lambda-2)(\lambda-3)(\lambda+6)=-C_3(\lambda^4-25\lambda^2+60\lambda-36)\). With \(C_3=1\):
\(C=(25,36,1)\) and \(T=60\). The network
\[
X_1\xrightarrow{\,1/25\,}2X_1+X_3,\qquad 2X_1+X_2\xrightarrow{\,1/22500\,}X_1+2X_2,\qquad
2X_2+X_3\xrightarrow{\,1/1296\,}X_2
\]
has, on the class \(x_1+x_2-x_3=60\), exactly the three positive equilibria \((25,36,1)\), \((50,18,8)\) and
\((75,12,27)\). The Jacobian restricted to \(\im\Gamma\) has determinant \(7/450\), \(-1/225\) and \(1/150\)
respectively, with eigenvalues \(\{-1.053,-0.0148\}\), \(\{-0.412,+0.0108\}\) and \(\{-0.384,-0.0174\}\): two stable
equilibria and a saddle.

\section{Computer verification}

The proof above needs no computer. As an independent consistency check, \texttt{scripts/certify\_family.py} builds
\(\Net_n\) for \(n=3,\dots,12\) and verifies in exact rational arithmetic (python-flint): the rank and kernels of
\(\Gamma\); the exponent formula; the strict sign alternation of the generalized-Vandermonde kernel vector for
\(\lambda_r=r\); the positivity of \(C\), \(T\) and \(\kappa\); that the \(n\) points \(x(r)\) satisfy \(\Gamma v(x)=0\) and
\(\varepsilon\cdot x=T\) exactly; that the restricted Jacobian has nonzero determinant at each of them; and that \(F\)
has exactly \(n\) sign changes (301 checks, all pass). \texttt{scripts/example\_n3.py} checks the instance of the
previous section with sympy.

\paperreferences
\end{document}
